\title{Direct Preference Optimization: Your Language Model is Secretly a Reward Model}

\begin{abstract}
Although large-scale unsupervised language models (LMs) acquire extensive general knowledge and some reasoning capabilities, exercising fine-grained control over their outputs remains challenging because their training is fully unsupervised. Common approaches to increasing the steerability of LMs involve obtaining human annotations regarding the relative merit of generated responses and then adjusting the unsupervised LM through fine-tuning to reflect these preferences, frequently leveraging reinforcement learning from human feedback (RLHF). Nevertheless, RLHF is a complicated and often unstable two-stage process: it first involves learning a reward model to capture human preferences, followed by fine-tuning the large, unsupervised LM with reinforcement learning to maximize this learned reward while avoiding significant divergence from the original model. In this work, we propose a novel parameterization for the reward model used in RLHF that makes it possible to directly obtain the optimal policy in closed form. This innovation enables the RLHF objective to be addressed through a straightforward classification loss, resulting in our algorithm, \textit{Direct Preference Optimization} (DPO). DPO offers stability, strong empirical results, and is computationally efficient, removing the necessity for sampling from the LM during fine-tuning or engaging in extensive hyperparameter optimization. Through experimental evaluations, we demonstrate that DPO aligns LMs with human preferences as effectively as, or better than, existing techniques. In particular, DPO surpasses PPO-based RLHF in controlling the sentiment of generated text, and is on par with or outperforms current approaches in tasks such as summarization and single-turn dialogue, all while offering significantly reduced complexity in implementation and training.
\end{abstract}

\section{Introduction}
Large-scale unsupervised language models (LMs) trained on extensive corpora have demonstrated remarkable emergent abilities~\citep{chowdhery2022palm, brown2020language, touvron2023llama, bubeck2023sparks}. Nevertheless, the data used to train these models originate from humans whose objectives, expertise, and priorities differ greatly. Not all these characteristics are worth emulating; for instance, although it is useful for an AI coding assistant to \textit{recognize} frequent programming errors to effectively correct them, we prefer the model to emulate the (possibly scarce) high-level programming proficiency captured in its training set when actually generating code. In a similar vein, we may want our language model to be \textit{cognizant} of widespread misconceptions—such as one held by 50\% of people—yet we would not want the model to affirm this erroneous belief in half of its responses. Thus, specifying the model's \emph{intended outputs and conduct} from its vast repertoire of \textit{knowledge and competencies} is essential to ensure that AI systems are reliable, high-performing, and controllable \citep{ouyang2022training}. Though prevailing techniques commonly employ reinforcement learning (RL) to align LMs with human values, we demonstrate that the standard RL-based preference optimization can actually be replaced with a straightforward binary cross-entropy objective, thereby streamlining the entire preference learning process.

In general, current approaches encourage desirable behaviors in language models by leveraging carefully constructed datasets of human preferences, reflecting behaviors regarded as beneficial and safe. This process, called preference learning, typically follows the preliminary phase of massive unsupervised pre-training on text. While supervised fine-tuning on human-annotated high-quality responses offers a direct means to encode human preferences, the dominant and most effective strategy involves reinforcement learning from human (or AI) feedback (RLHF/RLAIF; \citep{christiano2017deep, bai2022constitutional}). In RLHF, a reward model is trained on preference data provided by humans, and subsequently, RL is utilized to adjust the model’s policy so that its outputs receive higher reward scores—while ensuring the policy remains close to its pre-trained counterpart. Although RLHF is highly successful, producing language models with sophisticated conversational and coding skills, its pipeline is much more intricate than that of supervised learning, as it requires the training of multiple LMs and frequent policy sampling during training, thus resulting in considerable computational demands.

This paper demonstrates a method for optimizing language models to align with human preferences without relying on explicit reward modeling or reinforcement learning frameworks. We introduce \textit{Direct Preference Optimization (DPO)}, an algorithm that achieves the same optimization objective as traditional RLHF approaches—namely, reward maximization subject to a KL-divergence regularization—but with a much simpler design and training procedure. At a high level, DPO increases the log-likelihood ratio in favor of preferred responses over less preferred ones, yet unlike naive approaches based purely on probability ratios, it incorporates a dynamic, sample-specific importance weighting that mitigates the degeneration of the model commonly observed otherwise. DPO, akin to other existing algorithms, assumes a theoretical preference structure (for example, as described by the Bradley-Terry model; \cite{bradley1952rankanalysis}) that quantifies the alignment between reward functions and collected preference data. In contrast to conventional techniques, which use this theoretical model to establish a preference loss for reward model training and subsequently optimize a separate policy against the learned rewards, DPO re-frames the preference loss directly in terms of the policy through a change of variables. This allows the optimization of a policy over a dataset of human-labeled preferences via a straightforward binary cross-entropy loss, yielding an optimal policy corresponding to a latent reward function derived from these preferences.

The central contribution of this work is Direct Preference Optimization (DPO), a straightforward method for training language models on human preference data without RL. Empirical results indicate that DPO matches or surpasses current methods, including RLHF with PPO, across various preference-driven tasks such as sentiment control, summarization, and conversational modeling, tested on models up to 6B parameters in size.

\section{Related Work}

Increasingly large self-supervised language models have demonstrated the ability to handle certain tasks either in a zero-shot fashion \citep{radford2019language} or using a few-shot prompt paradigm \citep{gpt3,megatron,chowdhery2022palm}. Nonetheless, their effectiveness on specific downstream applications and their responsiveness to user instructions are markedly enhanced by subsequent fine-tuning on datasets composed of instructions and corresponding human-generated completions \citep{mishra-etal-2022-cross,sanh2022multitask,chung2022scaling,thoppilan2022lamda}. This process, known as `instruction-tuning,' empowers large language models (LLMs) to generalize beyond the instructions encountered during training, resulting in improved general-purpose usability \citep{chung2022scaling}. While instruction tuning has achieved considerable progress, it is often more practical to gather \textit{relative} human evaluations of model responses than to assemble datasets of expert completions. Consequently, later research has explored fine-tuning LLMs using collections of human preference data, yielding advancements in fields such as translation \citep{kreutzer-etal-2018-reliability}, summarization \citep{stiennon2022learning,ziegler2020finetuning}, story generation \citep{ziegler2020finetuning}, and instruction following \citep{ouyang2022training,ramamurthy2023is}. These approaches first involve training a neural reward model to align with recorded human preferences—often under a preference framework like the Bradley-Terry model \citep{bradley1952rankanalysis}—before subsequently tuning the language model to maximize this learned reward via reinforcement learning techniques such as REINFORCE \citep{williams1992reinforce}, proximal policy optimization (PPO; \cite{schulman2017proximal}), or similar methods \citep{ramamurthy2023is}. Related research also takes advantage of LLMs that have already been fine-tuned with human feedback to create synthetic preference datasets focused on attributes like safety or harmlessness \citep{bai2022constitutional}, relying on limited human oversight through textual rubrics that guide model evaluations. These strategies illustrate the merging of two principal research threads: the use of reinforcement learning for a variety of language modeling objectives~\citep{Ranzato2015SequenceLT,paulus2018a,wu2018learning} and the broader area of learning from human preference signals \citep{christiano2017deep,kupcsik2018learning}. Despite the promise offered by preference-based methods, fine-tuning large language models with reinforcement learning remains a significant technical hurdle; the present work introduces a theoretically sound methodology for optimizing relative preferences without the need for RL.

Beyond linguistic applications, the problem of learning policies from preference data has been explored within both bandit and reinforcement learning paradigms, giving rise to several methodologies. Contextual bandit frameworks that leverage preferences or rankings over actions instead of explicit reward signals are referred to as contextual dueling bandits (CDB; \cite{yue2012karmed,dudik2015contextual}). In these scenarios where absolute reward information is unavailable, the theoretical focus in CDBs shifts to the concept of a \textit{von Neumann winner}: a policy expected to outperform \textit{every} other policy with at least a 50\% probability \citep{dudik2015contextual}. Unlike CDBs, where preferences are collected interactively and online, human preference learning typically relies on a finite and static collection of pre-annotated action pairs \citep{yan2022human}. In a similar vein, \textit{preference-based RL} (PbRL) utilizes binary preferences, originating from an \textit{unknown} “scoring” function, in lieu of explicit reward feedback \citep{BusaFekete2014,ruiz2023dueling}. A diverse array of PbRL methods have been introduced—including strategies that accommodate off-policy preference datasets—yet most approaches initially estimate the hidden scoring or reward function before policy optimization proceeds \citep{jain2013learning,BusaFekete2014,christiano2017deep,sadigh2017active,kupcsik2018learning}. By contrast, our approach involves a single-stage procedure that aims to directly optimize the policy so that it aligns with observed preferences.

\section{Preliminaries}\label{section:prelims}

We briefly summarize the RLHF pipeline as presented by \citeauthor{ziegler2020finetuning} and further expanded by subsequent work \citep{stiennon2022learning, bai2022training, ouyang2022training}. The standard RLHF workflow is organized into three primary stages: (1) supervised fine-tuning (SFT); (2) collection of preference data and reward model training; and (3) reinforcement learning-based policy optimization.

\textbf{SFT}: The initial stage in RLHF involves adapting a pre-trained language model to the downstream domain(s) of interest (such as dialogue or summarization) through supervised learning on curated high-quality datasets. This process yields a model $\pi^\text{SFT}$.

\textbf{Reward Modelling Phase}: During this phase, the fine-tuned model $\pi^\text{SFT}$ generates candidate responses to given prompts $x$, yielding pairs $(y_1, y_2)\sim \pi^\text{SFT}(y \mid x)$. Human evaluators are then presented with each pair and asked to indicate a preference, denoted as $y_w\succ y_l \mid x$, where $y_w$ is favored over $y_l$ for the prompt $x$. Preferences are presumed to be dictated by a hidden reward function $r^*(y, x)$, which remains inaccessible. To statistically model these preferences, frameworks such as the Bradley-Terry (BT) model \cite{bradley1952rankanalysis} are commonly used (although more general models such as Plackett-Luce \citep{plackett1975analysis, luce2012individual} may be appropriate when rankings over more than two responses are available). The BT model specifies that the true human preference probability $p^*$ can be formalized as follows:

\begin{equation}\label{eq:bradley-terry}
    p^*(y_1\succ y_2 \mid x)=\frac{\exp\left(r^*(x, y_1)\right)}{\exp\left(r^*(x, y_1)\right) + \exp\left(r^*(x, y_2)\right)}.
\end{equation}

Suppose we have access to a fixed set of preference data $\mathcal{D}=\bigl\{x^{(i)}, y_w^{(i)}, y_l^{(i)}\bigr\}_{i=1}^N$ sampled from $p^*$. In this context, one can introduce a parameterized reward function $r_{\phi}(x, y)$, optimizing its parameters by maximizing the likelihood of observed preferences. By treating this as a binary classification task, the objective is defined via the negative log-likelihood:

\begin{equation}\label{eq:reward_model}
    \mathcal{L}_R(r_{\phi}, \mathcal{D}) = -\mathbb{E}_{(x, y_w, y_l)\sim \mathcal{D}}\bigl[\log \sigma(r_{\phi}(x, y_w)- r_{\phi}(x, y_l))\bigr]
\end{equation}

Here, $\sigma$ denotes the logistic sigmoid function. For language modeling scenarios, $r_{\phi}(x, y)$ typically initializes from the SFT model $\pi^\text{SFT}(y \mid x)$; an additional linear projection atop the final transformer layer then outputs a scalar reward estimate, following \cite{ziegler2020finetuning}. To encourage reward values with reduced variance, earlier works rescale the rewards such that $\mathbb{E}_{x,y\sim \mathcal{D}}\left[r_\phi(x, y)\right] = 0$ for each $x$.

\textbf{RL Fine-Tuning Phase}: In the subsequent RL phase, the trained reward function supplies evaluative signals for the language model’s outputs. Adopting the formulation from \citep{jaques2017sequence, jaques2020human}, one considers the objective:

\begin{equation}\label{eq:RL}
\max_{\pi_{\theta}}  \mathbb{E}_{x\sim \mathcal{D}, y\sim \pi_{\theta}(y \mid x)}\bigl[r_{\phi}(x, y)\bigr] - \beta\mathbb{D}_{\textrm{KL}}\bigl[\pi_{\theta}(y\mid x)\mid \mid \pi_\text{ref}(y\mid x)\bigr],
\end{equation}

Here, $\beta$ regulates the extent of divergence permitted from the base policy $\pi_\text{ref}$ (i.e., the original SFT model $\pi^\text{SFT}$). In practical implementations, the initial parameters of the language model policy $\pi_\theta$ are set to those of $\pi^\text{SFT}$. The KL divergence term acts as a regularizer, constraining the learned policy to remain close to distributions where the reward model is accurate, thus preserving diversity in outputs and avoiding collapse to narrow high-reward regions. Because sampling language is inherently discrete and non-differentiable, reinforcement learning approaches are required for optimization. The typical method \citep{ziegler2020finetuning, stiennon2022learning, bai2022training, ouyang2022training} constructs a modified reward $r(x, y) = r_{\phi}(x, y) -\beta (\log \pi_{\theta}(y\mid x) - \log \pi_\text{ref}(y\mid x))$, which is then maximized using Proximal Policy Optimization (PPO)~\cite{schulman2017proximal}.

\section{Direct Preference Optimization}\label{sec:DPO}

Considering the substantial challenges faced when scaling reinforcement learning for language model fine-tuning, we seek a streamlined method for optimizing policies based directly on preference data. Distinct from standard RLHF approaches that first infer a reward model before conducting RL-based optimization, our method exploits a particular parameterization of the reward function which allows its corresponding optimal policy to be written in closed form. As detailed in the following, our principal contribution lies in utilizing an analytic relationship mapping reward functions to their associated optimal policies, thus enabling transformation of a loss over reward models into a loss directly over policy parameters. This change-of-variable perspective allows policy learning to proceed without explicit, standalone reward model fitting, yet still conforms to preference models such as the Bradley-Terry formulation. Effectively, this setup unifies the

\textbf{Deriving the DPO objective.} Our starting point is the standard RL objective used in prior research, Eq.~\ref{eq:RL}, with a general reward function $r$. As established in earlier works~\citep{peters2007reinforcement, peng2019advantage, korbak2022reinforcement, go2023aligning}, it can be readily shown that the solution optimizing the reward while enforcing a KL constraint, as presented in Eq.~\ref{eq:RL}, is given by:

\begin{equation}\label{eq:op_policy}
    \pi_r(y\mid x) = \frac{1}{Z(x)}\pi_\text{ref}(y\mid x)\exp\left(\frac{1}{\beta}r(x, y)\right),
\end{equation}

where $Z(x) =\sum_{y}\pi_\text{ref}(y\mid x)\exp\left(\frac{1}{\beta}r(x, y)\right)$ serves as the normalization constant, also known as the partition function. A more thorough exposition can be found in Appendix \ref{app:derivation1}. Direct computation of $Z(x)$, even when employing an MLE estimate $r_{\phi}$ for the underlying true reward $r^*$, remains computationally demanding~\citep{korbak2022reinforcement, go2023aligning}, rendering this form less practical. Instead, Eq.~\ref{eq:op_policy} may be reformulated so that the reward is defined via the optimal policy $\pi_r$, the baseline policy $\pi_\text{ref}$, and the as-yet-unknown partition function $Z(\cdot)$. Specifically, applying the logarithm to both sides of Eq.~\ref{eq:op_policy} and rearranging terms yields:

\begin{equation}\label{eq:main_eq}
    r(x,y) =\beta \log \frac{\pi_r(y\mid x)}{\pi_\text{ref}(y\mid x)} + \beta \log Z(x).
\end{equation}

This change of variables can be equally well applied to the true reward $r^*$ and its optimal policy $\pi^*$. Importantly, the Bradley-Terry model is sensitive only to the reward difference between two outcomes, i.e., ${p^*(y_1 \succ y_2 \mid x) = \sigma(r^*(x, y_1) - r^*(x, y_2))}$. By substituting the expression from Eq.~\ref{eq:main_eq} for $r^*(x,y)$ into the preference model Eq.~\ref{eq:bradley-terry}, the partition function components eliminate each other, allowing the probability of human preference to be written exclusively in terms of $\pi^*$ and $\pi_\text{ref}$. Consequently, the optimal policy $\pi^*$ for RLHF under the Bradley-Terry preference framework is characterized as follows:

\begin{equation}\label{eq:objective}
    p^*(y_1\succ y_2 \mid x)=\frac{1}{1 + \exp\left(\beta \log \frac{\pi^*(y_2\mid x)}{\pi_\text{ref}(y_2\mid x)} - \beta \log \frac{\pi^*(y_1\mid x)}{\pi_\text{ref}(y_1\mid x)}\right)}
\end{equation}

The full derivation is given in Appendix~\ref{app:derivation2}. Although Eq.~\ref{eq:objective} is derived for the Bradley-Terry model, an analogous development holds for more general preference structures such as the Plackett-Luce models~\citep{plackett1975analysis, luce2012individual}, as described in Appendix~\ref{app:plackett_luce_models}.

Since we now have the human preference likelihood expressed in terms of the optimal policy (rather than via the reward function), we can construct a maximum likelihood training objective for a parametric policy $\pi_\theta$. Mirroring the formulation for reward modeling (see Eq.~\ref{eq:reward_model}), this gives the following policy training loss:

\begin{equation}\label{eq:optimum_model}
    \mathcal{L}_\text{DPO}(\pi_{\theta}; \pi_\text{ref}) = -\mathbb{E}_{(x, y_w, y_l)\sim \mathcal{D}}\left[\log \sigma \left(\beta \log \frac{\pi_{\theta}(y_w\mid x)}{\pi_\text{ref}(y

In this framework, we implicitly fit a reward function through an alternative form of parameterization, with the resulting optimal policy directly corresponding to $\pi_\theta$. Furthermore, because this approach is mathematically analogous to estimating a reparameterized Bradley-Terry model, it inherits desirable theoretical properties, including consistency when certain assumptions about the preference data are met \cite{bong2022generalized}. We explore further theoretical aspects of DPO and compare them to related literature in Section~\ref{sec:theory}.

\textbf{Mechanics of the DPO update.} To develop an intuition for how DPO operates, it is instructive to examine the gradient of its loss, $\mathcal{L}_\text{DPO}$. Differentiating with respect to the model parameters $\theta$ yields:

\begin{multline*}\label{eq:gradient}
    \nabla_\theta \mathcal{L}_\text{DPO}(\pi_\theta;\pi_\text{ref}) = \\ -\beta\mathbb{E}_{(x, y_w, y_l) \sim \mathcal{D}} \bigg[\underbrace{\sigma(\hat{r}_\theta(x, y_l) - \hat{r}_\theta (x, y_w))}_\text{greater emphasis if reward ranks are incorrect}\bigg[\underbrace{\nabla_\theta\log \pi(y_w \mid x)}_\text{encourage $y_w$} - \underbrace{\nabla_\theta\log\pi(y_l \mid x)}_\text{discourage $y_l$}\bigg]\bigg],
\end{multline*}

Here, the reward is given by $\hat{r}_\theta(x, y) = \beta \log \frac{\pi_\theta(y \mid x)}{\pi_\text{ref}(y \mid x)}$, which is implicitly defined using both the model $\pi_\theta$ and the reference policy $\pi_\text{ref}$ (additional details are provided in Section~\ref{sec:theory}). Intuitively, the loss gradient for $\mathcal{L}_\text{DPO}$ acts to make the model assign greater likelihood to preferred completions $y_w$ and lower likelihood to the less preferred $y_l$. A crucial detail is that each sample’s contribution is weighted by the extent to which the implicit reward $\hat{r}_\theta$ overvalues the disfavored output, modulated by the parameter $\beta$—effectively reflecting how much the reward misorders the examples, and the influence of the KL constraint. Empirically, our results highlight the necessity of this reweighting; omitting the weighting factor can cause the language model to collapse (see Appendix Table~\ref{tab:unlikelihood_generations}).

\textbf{DPO Overview.} The standard procedure for DPO can be summarized as follows: First, for each prompt $x$, generate candidate completions $y_1, y_2$ by sampling from $\pi_\text{ref}(\cdot \mid x)$. These completions are then annotated with human preference labels, resulting in an offline preference dataset $\mathcal{D} = \{x^{(i)}, y_w^{(i)}, y_l^{(i)}\}_{i=1}^N$. Second, the language model $\pi_\theta$ is trained to minimize the loss function $\mathcal{L}_\text{DPO}$, conditioned on the specified reference model $\pi_\text{ref}$, dataset $\mathcal{D}$, and a selected $\beta$. 

In practical applications, instead of generating new samples and collecting corresponding human annotations, practitioners often leverage publicly available preference datasets. When such datasets are created using $\pi^\text{SFT}$ as the sampling policy, we set $\pi_\text{ref} = \pi^\text{SFT}$ if possible. In the absence of an available $\pi^\text{SFT}$, $\pi_\text{ref}$ is instead obtained by fitting it to the preferred completions $(x, y_w)$ via maximum likelihood, i.e., ${\pi_\text{ref} = \argmax_{\pi}\mathbb{E}_{x, y_w \sim \mathcal{D}}\left[\log \pi(y_w \mid x)\right]}$. This initialization helps to reduce the discrepancy between the true (unobserved) reference distribution and the actual $\pi_\text{ref}$ used in DPO. Further specifics regarding experimental setup and hyperparameter choices are discussed in Appendix~\ref{app:implementation}.

\section{Theoretical Analysis of DPO} Here, we present a deeper theoretical perspective on the DPO framework, elaborating on its foundation, and drawing connections to challenges encountered in actor-critic-based RLHF algorithms, such as PPO~\cite{schulman2017proximal}, thus highlighting DPO’s advantages.

\label{sec:theory}

\subsection{Your Language Model Is Secretly a Reward Model} DPO circumvents the explicit fitting of a reward model as well as the use of reinforcement learning for policy optimization by employing a single maximum likelihood learning criterion. The loss in Eq. \ref{eq:main_eq} aligns with a Bradley-Terry model, where the reward is parameterized as $r^*(x, y) = \beta \log\frac{\pi^*_\theta(y \mid x)}{\pi_\text{ref}(y \mid x)}$. Training the parametric model $\pi_{\theta}$, in this context, can be seen as equivalent to optimizing the reward model in Eq. \ref{eq:reward_model} with an appropriate variable substitution. In what follows, we provide theoretical foundations for this reparameterization, demonstrate that it does not restrict the set of representable reward models, and prove that it can retrieve the optimal policy exactly. To begin, we introduce an equivalence relation among reward functions.

\begin{definition}
We define two reward functions $r(x, y)$ and $r'(x, y)$ as equivalent if and only if there exists a function $f$ such that $r(x, y)-r'(x, y) = f(x)$.     
\end{definition}

This notion constitutes an equivalence relation, effectively dividing the collection of all reward functions into disjoint equivalence classes. We establish two key lemmas as follows:

\begin{lemma}\label{lemma:same_prefrence} Within the framework of Plackett-Luce, and specifically within Bradley-Terry models, any pair of reward functions belonging to the same equivalence class gives rise to an identical preference distribution.
\end{lemma}

\begin{lemma}\label{lemma:same_policy}
    Any two reward functions that reside in the same equivalence class lead to the same optimal policy when considered under the constrained RL formulation.
\end{lemma}

Since their proofs are elementary, we postpone their details to Appendix \ref{app:lemma1}. The first lemma underscores a familiar problem of under-specification in the Plackett-Luce family \cite{plackett1975analysis}: identifiability issues compel us to impose supplementary constraints in order to ensure the uniqueness of the MLE solutions for Eq. \ref{eq:reward_model} \cite{bong2022generalized}. The second lemma clarifies that every member of a reward class determines the same optimal policy, which implies that for purposes of policy recovery, it suffices to identify any representative of the optimal reward class. In Appendix~\ref{app:thm1}, we provide a proof for the following theorem:

\begin{theorem}\label{thm:main}
    Assuming standard regularity conditions, every reward equivalence class compatible with the Plackett-Luce model (and the Bradley-Terry model in particular) admits the reparameterization ${r(x, y) = \beta \log \frac{\pi(y\mid x)}{\pi_\text{ref}(y\mid x)}}$ for some policy $\pi(y\mid x)$ and a fixed reference policy $\pi_\text{ref}(y \mid x)$.
\end{theorem}

\begin{sproof}
    Take an arbitrary reward function $r(x, y)$, which defines a unique optimal policy $\pi_r(y \mid x)$ according to Eq. \ref{eq:op_policy}. We show that there exists a representative in the equivalence class of $r$ that is expressible via the proposed reparameterization. We introduce the transformation $f$ as  
\begin{equation}
    f(r; \pi_\text{ref}, \beta)(x, y) = r(x, y) - \beta\log\sum_{y}\pi_\text{ref}(y\mid x)\exp\left(\frac{1}{\beta}r(x, y)\right)
\end{equation}
This mapping $f$ normalizes the reward function by subtracting the log-partition function (computed under $\pi_\text{ref}$ and $r$), which only depends on $x$. Therefore, $f(r; \pi_\text{ref}, \beta)(x, y)$ lies within the equivalence class of $r(x, y)$. Substituting the right-hand side of Eq.~\ref{eq:main_eq}—which is valid for all reward functions—yields $f(r; \pi_\text{ref}, \beta)(x, y) = \beta \log \frac{\pi_r(y\mid x)}{\pi_\text{ref}(y\mid x)}$. Hence, this projection produces a reward function in the same equivalence class, adopting the desired structure, and demonstrates the generality of the suggested reparameterization for our reward model.
\end{sproof}

An alternative interpretation of Theorem~\ref{thm:main} is that it characterizes the precise choice of reward function, within each equivalence class, that is singled out by the DPO reparameterization; specifically, it identifies the reward function that fulfills the following condition:

\begin{equation}\label{eq:lag_p}
     \sum_{y}\underbrace{\pi_\text{ref}(y\mid x)\exp\left(\frac{1}{\beta}r(x, y)\right)}_{=\pi(y\mid x)\text{, using Thm.~\ref{thm:main} reparam.}} = 1,
\end{equation}

which ensures that $\pi(y\mid x)$ constitutes a proper probability distribution (all probabilities are non-negative and sum to unity). Referring back to Eq.~\ref{eq:op_policy}, it is apparent that Eq.~\ref{eq:lag_p} corresponds to the partition function for the optimal policy associated with the reward $r(x, y)$. The central contribution of the DPO framework lies in demonstrating that by imposing particular constraints on the otherwise underdetermined Plackett-Luce (or specifically, Bradley-Terry) family of preference models, one can maintain the generality of the reward model class, while rendering the partition function in Eq.~\ref{eq:op_policy} (and hence the optimal policy) computationally manageable for arbitrary prompts $x$.

\subsection{Instability of Actor-Critic Algorithms} This formalism also facilitates the analysis of instability issues encountered in conventional actor-critic approaches—such as PPO—used in RLHF pipelines. Here, our focus is on the reinforcement learning fine-tuning stage as described in Section~\ref{section:prelims}. Our methodology draws parallels with the control as inference paradigm \cite{levine2018reinforcement} in the context of the constrained RL setup specified in \ref{eq:RL}. Let us consider a parameterized family of policies $\pi_{\theta}(y\mid x)$ and seek to minimize the divergence $\mathbb{D}_{\text{KL}}[\pi_{\theta}(y|x) \mid \mid \pi^*(y\mid x)]$, where $\pi^*$ represents the optimal policy, defined by the reward $r_{\phi}(y, x)$ via Eq.~\ref{eq:optimum_model}. A straightforward algebraic manipulation yields the following objective:

\begin{equation}\label{eq:AC}
    \max_{\pi_{\theta}}\mathbb{E}_{\pi_{\theta}(y\mid x)}\bigg[\underbrace{r_{\phi}(x, y) -\beta\log\sum_{y}\pi_\text{ref}(y\mid x)\exp\left(\frac{1}{\beta}r_{\phi}(x, y)\right)}_{f(r_{\phi}, \pi_\text{ref}, \beta)} - \underbrace{\beta\log\frac{\pi_{\theta}(y\mid x)}{\pi_\text{ref}(y\mid x)}}_{\text{KL}}\bigg]
\end{equation}

The objective here is identical to that targeted in earlier studies 
\citep{ziegler2020finetuning, stiennon2022learning, bai2022training, ouyang2022training}, employing the DPO-equivalent reward structure associated with the reward function $r_{\phi}$. Within this framework, the normalization component in $f(r_{\phi}, \pi_\text{ref}, \beta)$ can be viewed as the soft value function corresponding to the reference policy $\pi_\text{ref}$. Although this normalization term is not necessary to achieve optimality, omitting it can lead to high-variance policy gradients, thereby causing instability during learning. This issue can be mitigated by incorporating a value function estimator, though optimizing such an estimator often proves challenging. As an alternative, previous research has utilized a human completion as a baseline for normalization, which acts as a one-sample Monte Carlo approximation of the normalization factor. The DPO reparameterization, however, produces a reward function that is inherently baseline-free.

\section{Experiments}
In the following section, we provide an empirical analysis of DPO’s capacity to learn policies directly from preference data. We begin with a controlled text-generation environment, investigating how well DPO balances reward maximization and KL-divergence minimization with respect to the reference policy, compared to established algorithms like PPO. We subsequently extend our evaluation to encompass larger-scale models and more complex RLHF tasks, such as dialogue and summarization. Our results indicate that, with minimal hyperparameter adjustment, DPO generally matches or surpasses strong baselines—including RLHF with PPO and the selection of the best out of $N$ sampled trajectories using a trained reward model. Prior to presenting these empirical findings, we outline our experimental methodology; comprehensive experimental details are provided in Appendix~\ref{app:exp_details}.

\textbf{Tasks.} The experimental setup involves three open-ended text generation tasks. Across all experiments, each method derives a policy based on a preference dataset $\mathcal{D}=\bigl\{x^{(i)}, y_w^{(i)}, y_l^{(i)}\bigr\}_{i=1}^N$. For the task of \textbf{controlled sentiment generation}, $x$ represents the initial segment of a movie review sourced from the IMDb dataset \cite{maas-EtAl:2011:ACL-HLT2011}; the objective for the policy is to produce $y$ with positive sentiment. To systematically assess performance, we create preference pairs over outputs using an existing sentiment classifier, selecting those for which $p(\text{positive}\mid x,y_w)>p(\text{positive}\mid x,y_l)$. For SFT, GPT-2-large is fine-tuned on the training portion of the IMDb reviews until convergence (additional information in App~\ref{app:sentiment_details}). In the \textbf{summarization} task, $x$ consists of a Reddit forum post, and the model must generate a summary $y$ capturing the main content. Building on existing literature, this work adopts the Reddit TL;DR summarization dataset \citep{volske-etal-2017-tl} supplemented by human preferences curated by \citeauthor{stiennon2022learning}. We use an SFT checkpoint that was previously fine-tuned on human-authored forum post summaries\footnote{\url{https://huggingface.co/CarperAI/openai_summarize_tldr_sft}} within the TRLX \citep{leandro_von_werra_2023_7790115} RLHF infrastructure. The associated human preference labels originate from \citeauthor{stiennon2022learning} and pertain to outputs from a similarly-trained, but distinct, SFT model. Lastly, in the \textbf{single-turn dialogue} scenario, $x$ corresponds to a user-initiated prompt ranging from scientific inquiries to personal advice requests. Here, the policy is expected to provide a response $y$ that is both useful and engaging. We rely on the Anthropic Helpful and Harmless dialogue dataset \citep{bai2022training}, which includes 170,000 human-assistant conversations. Each dialogue is concluded by two alternative model-generated responses and an annotation specifying which was favored by human raters. Unlike prior tasks, no pre-trained SFT model exists for this data, prompting us to fine-tune a general-purpose language model solely on the preferred responses to obtain the SFT baseline.

\textbf{Evaluation.} Our experimental evaluation employs two distinct methodologies. To examine how well each algorithm optimizes the objective of maximizing reward under constraints, we assess their performance in controlled sentiment generation by plotting the trade-off frontier between reward achieved and KL-divergence from the original policy; this is feasible due to access to an explicit, ground-truth reward function in the form of a sentiment classifier. In contrast, because real-world settings lack a ground truth reward function, we measure algorithm performance through their \textit{win rate} relative to a baseline policy. Here, GPT-4 serves as a stand-in for human judgment, evaluating summary quality and response helpfulness in summarization and single-turn dialogue, respectively. For the summarization task, baseline comparisons are made using reference summaries from the test split; for dialogue, the baseline is the preferred response designated in the dataset. While prior research indicates that language models can function as more reliable automated evaluators than existing automated metrics \citep{Chen2023ExploringTU}, we also conduct a human study, described in Sec.~\ref{sec:human-judgments}, to validate the appropriateness of using GPT-4 for our evaluations. We observe that GPT-4’s evaluations are closely aligned with human judgments, with agreement rates between humans and GPT-4 on par with or surpassing the consistency seen between human annotators.

\textbf{Methods.} Alongside DPO, our evaluation includes a range of established techniques for training language models according to human preferences. As the simplest baselines, we use zero-shot prompting with \textbf{GPT-J} \citep{gpt-j} for summarization and 2-shot prompting with \textbf{Pythia-2.8B} \citep{biderman2023pythia} for the dialogue setting. Further, we evaluate the \textbf{SFT} model as well as \textbf{Preferred-FT}, which involves supervised fine-tuning on the preferred output $y_w$, sourced either from the SFT model (in sentiment control and summarization tasks) or from a general language model (for single-turn dialogue). Another related technique is the \textbf{Unlikelihood} approach~\citep{welleck2019neural}, where training maximizes the probability of $y_w$ and concurrently reduces the probability assigned to $y_l$, optionally scaling the unlikelihood component by a parameter $\alpha\in[0,1]$. We also include \textbf{PPO} \citep{schulman2017proximal} using a reward model learned from preference annotations, as well as \textbf{PPO-GT}, an oracle method that utilizes the true reward function available in the controlled sentiment setting. Within our sentiment analysis experiments, we adopt two PPO-GT variants: a standard implementation \cite{leandro_von_werra_2023_7790115}, and a modified version with reward normalization and further hyperparameter adjustments to enhance results (similar enhancements are also applied when deploying PPO with learned rewards). Lastly, we compare against the \textbf{Best of $N$} method, where $N$ completions are sampled from the SFT model (or Preferred-FT for dialogue), and the completion with the highest score under a learned reward function is chosen. While this approach can yield strong results by decoupling reward model quality from PPO optimization, it incurs significant computational cost as it necessitates generating $N$ outputs per input at inference time, which becomes infeasible for even moderately large $N$.

\subsection{How well can DPO optimize the RLHF objective?}

The reward maximization objective in standard RLHF approaches imposes a KL-divergence constraint, striking a balance between maximizing reward and keeping the learned policy close to a reference policy. Hence, when evaluating different methods, it is important to consider both the attained reward and the extent of KL-divergence; simply attaining marginally higher rewards at the cost of significantly increased KL divergence is often not advantageous. Figure~\ref{fig:frontier-tldr-main} presents the trade-off curves between reward and KL for several algorithms on the sentiment task. For each method, we conduct several training experiments with varying degrees of policy conservativeness by altering the relevant hyperparameter (target KL $\in\{3,6,9,12\}$ for PPO, $\beta \in \{0.05,0.1,1,5\}$, $\alpha\in\{0.05,0.1,0.5,1\}$ for unlikelihood, and using different random seeds for preferred-FT), amounting to a total of 22 separate runs. Every 100 training steps up to convergence, policies are evaluated against a fixed set of test prompts, and we compute both the average reward using the true reward function and the mean sequence-level KL\footnote{Specifically, this is the aggregate sum of KL-divergences calculated at each timestep.} relative to the reference policy $\text{KL}\left(\pi\mid \mid \pi_\text{ref}\right)$. Our results reveal that DPO delineates the most favorable frontier—achieving superior rewards with consistently low KL values. This finding is especially significant for several reasons. Notably, while both DPO and PPO are designed to optimize the same objective, DPO achieves this with much greater efficiency, yielding a reward/KL profile that uniformly outperforms PPO. Additionally, DPO establishes a superior reward-KL curve even when PPO is provided access to exact, ground truth rewards (PPO-GT).

\subsection{Can DPO scale to real preference datasets?}
\label{sec:dpo-real-datasets}
We proceed to investigate the fine-tuning capabilities of DPO on tasks involving summarization and single-turn dialogue. In the context of summarization, it is well established that automatic metrics such as ROUGE show weak alignment with human judgment~\citep{stiennon2022learning}, and previous research demonstrates that optimizing LMs using PPO on human preferences can produce summaries that are better received by users. Our evaluation consists of generating completions for the test split of the TL;DR summarization dataset and calculating the mean win rate in comparison to the reference completions. Sampling for each approach is performed across a range of temperatures from 0.0 to 1.0, with results displayed in Figure~\ref{fig:frontier-tldr-main} (right). For all experiments, DPO, PPO, and Preferred-FT are applied to the same underlying GPT-J SFT model\footnote{\url{https://huggingface.co/CarperAI/openai_summarize_tldr_sft}}. Our results indicate that DPO attains a win rate of roughly 61\% at a temperature of 0.0, surpassing PPO, which achieves approximately 57\% under its optimal sampling condition. DPO also secures a higher peak win rate than the best of $N$ baseline approach. It is worth mentioning that no extensive tuning of DPO's $\beta$ hyperparameter was conducted, implying that the reported figures could be conservative estimates of DPO's true performance. Furthermore, DPO demonstrates enhanced resilience to variations in sampling temperature relative to PPO, which can deteriorate to the level of the base GPT-J model at higher temperatures. In contrast, Preferred-FT offers little to no improvement over the initial SFT model. Additional direct comparison through human assessments is presented in Section~\ref{sec:human-judgments}, revealing that DPO samples generated at a temperature of 0.25 were preferred 58\% of the time over those produced by PPO at the same temperature.

For single-turn dialogue tasks, we assess the various approaches using the subset of the test split from the Anthropic HH dataset \citep{bai2022training} that consists of a single human-assistant exchange. In the case of GPT-4 evaluations, we utilize the preferred responses in the test set as the reference standard, measuring each method's win rate accordingly. Due to the absence of an established SFT baseline, our experiments begin with a pre-trained Pythia-2.8B model. We apply Preferred-FT to train a reference model using selected completions, ensuring the generated completions align with the model's data distribution, followed by DPO training. In addition, we benchmark against the strongest candidate from 128 Preferred-FT completions (as analysis revealed that the Best of $N$ baseline reaches its maximum at 128 for this task; refer to Appendix Figure~\ref{fig:best-of-n}) and a 2-shot prompting configuration using the Pythia-2.8B base model. Our findings indicate that DPO matches or surpasses these baselines, particularly at optimal temperatures. We further include an RLHF system trained via PPO on the Anthropic HH dataset \footnote{\url{https://huggingface.co/reciprocate/ppo_hh_pythia-6B}}, developed by a recognized group \footnote{\url{https://github.com/CarperAI/trlx/tree/main/examples/hh}}. Despite this, we could not identify a prompt or sampling temperature for PPO that yields better performance than the underlying Pythia-2.8B. Drawing on evidence from TL;DR and recognizing that Best of 128 and PPO both maximize the same reward objective, we interpret Best of 128 as a practical approximation of PPO-level results. Taken together, DPO distinguishes itself as the sole computationally efficient strategy that surpasses the reference completions within the Anthropic HH dataset, and its outcomes are comparable to or exceed those of the compute-intensive Best of 128 benchmark. Lastly, Figure~\ref{fig:dialogue-main} demonstrates that DPO rapidly attains peak performance.

\subsection{Generalization to a new input distribution}

To extend our comparison of PPO and DPO under distributional changes, we tested the policies obtained from our Reddit TL;DR summarization setup on a distinct dataset: the news article summaries from the CNN/DailyMail test split \citep{nallapati-etal-2016-abstractive}. We applied the optimal sampling temperatures previously identified for TL;DR (0 and 0.25). Table~\ref{tab:ood} presents these findings. For evaluation, we calculated the win rate of GPT-4 when pitted against reference summaries from each dataset, utilizing the same GPT-4 (C) evaluation prompt as in the TL;DR experiment, with "forum post" replaced by "news article". DPO maintained a clear advantage over the PPO policy on this novel input domain. These results preliminarily indicate that DPO achieves a degree of generalization on par with PPO, despite not leveraging the supplementary unlabeled Reddit TL;DR prompts available to PPO.

\subsection{Validating GPT-4 judgments with human judgments}
\label{sec:human-judgments}
We ran a human evaluation to assess how reliably GPT-4 reflects human preferences, employing the outcomes from the TL;DR summarization study as well as two alternative GPT-4 prompts. The \textbf{GPT-4 (S)} prompt takes a straightforward approach, asking which summary more effectively conveys the key information in the post. By contrast, the \textbf{GPT-4 (C)} prompt additionally inquires which summary is more concise. The latter was examined since GPT-4 tends to select longer, more repetitive outputs under the former prompt compared to human raters. Complete versions of the prompts are available in Appendix~\ref{app:prompts}. For evaluation, we selected three comparison points: the top-performing method (DPO, temp. 0.25), the lowest-performing method (PPO, temp. 1.0), and an intermediate baseline (SFT, temp. 0.25), all pitted against PPO under its best greedy temperature. This approach was intended to ensure a broad representation of summary quality; for the DPO and PPO-1 comparisons, multiple human ratings were gathered given the limited pool of human annotators. In both prompt conditions, GPT-4’s judgments exhibited agreement with humans at levels comparable to inter-annotator agreement, supporting GPT-4’s role as an effective surrogate for human assessment. Notably, the \textbf{GPT-4 (C)} prompt yielded win rates that better matched those of human raters, motivating its use for the principal results in Section~\ref{sec:dpo-real-datasets}. Details regarding the human evaluation setup, such as the rater web interface and information about volunteer annotators, are included in Appendix~\ref{app:human-study}.

\section{Discussion}
Preference-based learning provides a robust and extensible approach for developing competent, value-aligned language models. In this work, we have presented DPO, a straightforward methodology for learning from preference data without relying on reinforcement learning methods. Instead of reformulating the preference learning objective to conform to traditional RL frameworks in order to leverage existing RL tools, DPO establishes a direct correspondence between language model policies and reward mechanisms, making it possible to train language models according to human preferences \textit{directly} using a standard cross-entropy loss—eschewing reinforcement learning techniques and without sacrificing generality. Remarkably, DPO attains performance on par with or surpassing established RLHF methods, such as those utilizing PPO, all while requiring minimal hyperparameter adjustment; this effectively lowers the obstacles associated with training language models guided by human preferences.

\textbf{Limitations \& Future Work.} Our findings prompt a number of avenues for future research. For example, to what extent do DPO-trained policies extrapolate to out-of-distribution inputs compared with approaches that rely on explicit reward functions? Preliminary evidence points to DPO-based models exhibiting generalization capabilities comparable to those trained with PPO, yet further empirical investigation is warranted. Additional questions arise, such as whether incorporating self-labeling techniques within the DPO framework could allow for productive use of unlabeled prompts. Another area to investigate is how excessive reward maximization occurs in the direct preference optimization regime, and whether the slight reduction in performance observed in Figure~\ref{fig:dialogue-main}-right reflects such dynamics. While our evaluation focused on models up to 6B parameters, a promising trajectory lies in scaling DPO training to much larger, state-of-the-art language models. With respect to assessment methodologies, we observe that GPT-4-derived win rates are sensitive to the choice of prompt, inviting further work into optimizing automated evaluation strategies for reliability and consistency. Lastly, DPO's principles hold promise for broader application, potentially extending to the training of generative models in diverse domains beyond language models informed by human preferences.